\scanpage{217}
Lien avec la variété des modules pour des \unsure{multip.|multips|motifs} polarisées rigidifiées : Si \deleted{\unsure{$A$|$S$}} $\unsure{\mathcal M|\mathscr M}$ est une telle variété modulaire, alors\unsure{,} un morphisme $X\unsure{\mathrel{\text{---}}|\longrightarrow}\mathcal M$ est déterminé quand on connaît sa valeur \uline{en un point}, plus l'application $\pi_1(x)\mathrel{\text{---}}\pi_1(\mathcal M)$ ......

C'est une propriété intrinsèque importante de ces $\mathcal M$, et une propriété de cette nature mérite d'être étudiée pour \unsure{ses|les} propres mérites .... Peut-être \unsure{peuvent|pourront} elles \unsure{ttes|toutes} être \deleted{\unsure{considérées}} \inserted{interprétées} comme des \deleted{\unsure{schémas|espaces}} \inserted{variétés} modulaires pour certains types de motifs.

\begin{center}\rule{3cm}{0.5pt}\end{center}

Introduire la notion de schéma \inserted{motiviquement} \unsure{ponctué|pointé} \inserted{\tikz[baseline=(a.base)]\node[draw,ellipse,inner sep=1pt](a){(à val. dans un anneau $k$)};}, et morphismes de tels, qui définissent des morphismes correspondants pour les
\[
\begin{aligned}
&G(X,\unsure{\xi|\zeta})\mathrel{\text{---}}G(Y,\unsure{\eta|y}),\qquad\text{donc pour}\\
&G(X,\xi)(k)=\pi_1(X,\xi)\mathrel{\text{---}}G(Y,\eta)(k)=\pi_1(Y,\eta).
\end{aligned}
\]
% archive sheet number: 216

